\chapter{The prototype so far}\label{ch:prototype}
% Draft by Claude from our old iteration notes, edit freely.
% This chapter is a snapshot up to a milestone, not a list of iterations. Keep adding to it
% until the semester ends. For each part: what we did, and WHY.

This chapter describes the prototype as it is on \fillin{milestone date}. Instead of listing
every version, each section describes one part of the device: what it does now, and why we
made it that way. Details of each printed version (photos, notes) are kept in the
\texttt{iterations/} folder of the project repository.

\section{How the device works}\label{sec:how-it-works}

The device is a small, automated version of the Hirst trap described in
Chapter~\ref{ch:introduction}:

\begin{enumerate}
  \item A fan moves outdoor air onto a strip of clear sticky tape, so particles stick to it.
  \item A stepper motor moves the tape on, so the sampled piece lies over the camera.
  \item An LED shines through the tape, and the camera underneath takes a picture.
  \item The picture is saved with the time and all settings. Later it is sent to a server.
  \item This repeats. The used tape is kept on a roll, so it can be checked under a microscope later.
\end{enumerate}

A Raspberry Pi~3~B+ controls everything. \fillin{Photo of the current prototype.}

\section{Enclosure and tape transport}\label{sec:enclosure}

We started from the earlier design (Section~\ref{sec:research-question}), but drew our own 3D
model from scratch. The reasons were:

\begin{itemize}
  \item \textbf{Easy to build.} The goal is that anyone can build one, so we avoided screws
  and other extra parts where we could. Parts hold by friction or slot into each other.
  \item \textbf{One print.} We have access to a larger 3D printer, so parts can be bigger and
  fewer.
\end{itemize}

The main changes over our versions so far were:

\begin{itemize}
  \item \textbf{Tape roll holder.} The first holder had the wrong size. It now holds the roll
  and the tape stays on.
  \item \textbf{Motor mount.} The motor first sat loose. It now friction fits, with a small gap
  so gravity helps hold it. The mount is still a little too small, which makes the frame bend.
  \future{Fix the motor mount size.}
  \item \textbf{Camera slot and lid.} The first versions were made before we had the camera, so
  there was no proper hole for it. The camera now slots in exactly, and the lid has a hole for
  the ribbon cable. We first designed a complicated lid mechanism, then replaced it with a lid
  that is simply held by friction. Simple turned out to be better.
\end{itemize}

\paragraph{Still missing.}
\future{A mount for the fan.}
\future{A place for the LED inside the box, at a fixed distance above the tape.}
\future{A nozzle or slit that guides the air onto a small spot of the tape.}

\paragraph{Blowing or sucking.}
The reference trap sucks air through a narrow slit onto the tape. That gives a known airflow
and a sharp, narrow line of particles, which is needed to report pollen per cubic metre of air.
For the first prototype we chose to blow air onto the tape instead, because it is much simpler
to build. A printed nozzle between fan and tape is a possible middle ground.
\future{Measure the airflow, and decide whether to change to sucking through a slit.}

\section{Electronics}\label{sec:electronics}

\begin{itemize}
  \item \textbf{Stepper motor.} A 28BYJ-48 with a ULN2003 driver board moves the tape. It is
  slow and weak, but cheap and strong enough for tape.
  \item \textbf{Fan.} A 5~V fan, switched by a MOSFET, because a Raspberry Pi pin cannot power a
  fan directly. A diode protects the MOSFET.
  \item \textbf{Camera.} The Raspberry Pi HQ camera with the lens removed
  (Section~\ref{sec:imaging-tests}).
  \item \textbf{LED.} An RGB LED with a resistor for each colour, connected directly to three Pi pins.
\end{itemize}

\paragraph{Why an RGB LED.}
Our first LED was blue, and pure blue light gave poor pictures
(Section~\ref{sec:imaging-tests}). With an RGB LED the software can choose the colour and
brightness, so we can test which light works best.
